\subsubsection{Parameterzuordnung aus städtischem Baumkataster}

Die abschließende Stufe der Datenanreicherung fusioniert die detektierten Baumspitzen mit den administrativen Bestandsdaten des Leipziger Baumkatasters. Während die vorangegangenen Schritte die Bäume rein geometrisch und spektral beschreiben, dient dieser Abgleich der Identifikation der realen Bäume. Administrative Attribute wie Gattung, Pflanzjahr oder amtliche Baumnummer sollen in den LiDAR-Datensatz integriert werden. \cite{munzingerMappingUrbanForest2022}

Analog zu den vorherigen Prozessschritte erfolgt der Abgleich kachelweise. Zur Steigerung der Recheneffizienz lädt eine räumliche Index-Abfrage selektiv nur die Katastereinträge, die innerhalb der aktuellen Kachel zuzüglich eines Puffers von 20 m liegen. \cite{munzingerMappingUrbanForest2022}

Die Zuordnung der Datensätze basiert auf einer 3D-Distanzberechnung. Während das Baumkataster präzise X- und Y-Koordinaten bereitstellt, muss die Z-Komponente näherungsweise über Sachdaten abgeleitet werden. Die Baumhöhe ist in 165.756 von insgesamt 177.334 Katastereinträgen (ca. 93,5 ~\%) hinterlegt. Für die verbleibenden Daten wird die Höhe über den doppelten Kronendurchmesser geschätzt, da im urbanen Bestand eine allometrische Korrelation zwischen horizontaler Kronenausdehnung und vertikalem Wachstum vorliegt. Fehlt auch dieser Wert, wird der statistische Median aller Baumhöhen innerhalb des lokalen Abschnitts herangezogen. \cite{munzingerMappingUrbanForest2022}

Sobald für beide Datensätze 3D-Koordinaten vorliegen, wird die Zuweisung über einen adaptiven Suchraum räumlich eingegrenzt. Liegt der Kronendurchmesser im Kataster vor, entspricht der Suchradius dem halben Durchmesser zuzüglich einer Toleranz von 1,0 m. Anderenfalls skaliert der Radius dynamisch in Abhängigkeit der Baumhöhe innerhalb von 2,0 bis 8,0 m. Auf Basis dieser Radien wird eine paarweise Distanzmatrix für diese Radien aufgestellt. Baumspitzen außerhalb dieses Radius erhalten einen unendlichen Kostenwert, wodurch sie für eine Verknüpfung ausgeschlossen werden. \cite{murrayARBORNewFramework2019}

In dichten Beständen wie in Parkanlagen würde ein klassischer Nearest-Neighbor-Ansatz aufgrund von Kronenüberlagerungen zu fehlerhaften Zuordnungen führen. \cite{munzingerMappingUrbanForest2022} Deshalb wird die Fusion als globales Optimierungsproblem formuliert und über den Ungarischen Algorithmus für die gesamte Kachel gelöst. Das Verfahren minimiert die Gesamtsummer der Distanzen innerhalb der Kostenmatrix und garantiert eine eindeutige Zuordnung. \cite{murrayARBORNewFramework2019}

Nach erfolgreichem Abgleich wird das Baumsegment im GPKG-Datensatz als Match verifiziert und mit alle Attributen des Katasters angereichert.
